\chapter{Extensions natives}\label{ch:pl:native-extensions}

Un chercheur a réécrit en C++ l'étape interne d'un signal et l'a appelée depuis Python une fois par tick, un million de fois. Elle était plus lente
que la fonction Python qu'elle remplaçait~: 118~ns par appel contre 60~ns. Appelée une seule fois sur l'ensemble du tableau d'un million de valeurs, la
même boucle C++ a coûté 1.4~ns par élément, soit plus de deux fois plus vite que la routine de bibliothèque qu'elle devait battre. Tout
ce qui concerne une \omterm{def:pl:native-extensions:extension}{extension native} se joue à la frontière entre l'interpréteur et le code compilé~: à quelle fréquence on la franchit, ce qu'il faut
convertir au passage, et qui possède la mémoire de chaque côté. Ce chapitre lie un noyau écrit en C++20 et en Rust à Python
de deux façons, avec pybind11 et par une interface C chargée avec \texttt{ctypes}, fait passer les tableaux numpy sans les copier,
et mesure la frontière (\texttt{firm.natext}).

\section{Pourquoi et quand quitter Python}

\begin{definition}[Extension native, liaison de langage]\label{def:pl:native-extensions:extension}
Une \emph{extension native}\index{extension native} (native extension) est du code compilé (C, C++, Rust) chargé dans le processus Python et appelé comme une
fonction Python. Une \emph{liaison de langage}\index{liaison de langage} (language binding) est la couche qui la rend appelable~: elle convertit les arguments et
les résultats entre les représentations des deux langages, signale les erreurs à travers la frontière et précise qui possède chaque morceau de mémoire.
\end{definition}

Le chapitre~8 a donné l'ordre de préférence~: d'abord la \omterm{def:pl:python-at-scale:array}{programmation par tableaux}, qui repose déjà sur du code compilé écrit par d'autres~; ensuite une
\omterm{def:pl:native-extensions:extension}{extension native} pour ce que les tableaux ne savent pas exprimer --- une boucle dont l'état change à chaque élément (un carnet d'ordres, une position dans la file, une machine à
états), une fusion de deux suites triées, une récursion avec des branchements. Les deux noyaux du chapitre sont un exemple de chaque sorte~: une moyenne mobile exponentielle
(que numpy sait exprimer comme un filtre, de sorte que la version native rivalise avec une bibliothèque) et une recherche as-of de dates triées
en une seule passe de fusion (que numpy effectue avec une recherche dichotomique par élément, $O(n \log m)$ au lieu de $O(n + m)$).

\omcode{code/firm/natext/cpp/firm_natext.hpp}{9}{38}{Les deux noyaux en C++20, indépendants de Python~: la moyenne, une étape de celle-ci, et la
recherche as-of par une passe de fusion.}\label{lst:pl:native-extensions:kernels}

\section{La frontière~: appels, types et verrou de l'interpréteur}

\begin{definition}[Interface de fonctions étrangères, interface binaire d'application]\label{def:pl:native-extensions:ffi}
Une \emph{interface de fonctions étrangères}\index{interface de fonctions étrangères} (foreign function interface) permet à du code écrit dans un langage d'appeler des fonctions compilées à partir d'un autre.
Elle repose sur une \emph{interface binaire d'application}\index{interface binaire d'application} (application binary interface)~: la convention au niveau machine (la façon dont les arguments
passent dans les registres et sur la pile, la disposition des structures, les noms des symboles), à laquelle l'appelant et l'appelé, compilés
séparément, se conforment tous deux~; la convention C est celle que tous les langages savent parler.
\end{definition}

\begin{definition}[Coût de franchissement de la frontière]\label{def:pl:native-extensions:crossing}
Le \emph{coût de franchissement de la frontière}\index{coût de franchissement de la frontière} (boundary-crossing cost) d'un appel natif est le travail fixe effectué à chaque appel, indépendamment de
la taille de ses données~: trouver la fonction, vérifier et convertir chaque argument à partir d'un objet Python, relâcher puis reprendre le verrou de l'interpréteur,
reconvertir le résultat et gérer les erreurs.
\end{definition}

\begin{omfigure}
\begin{tikzpicture}[font=\scriptsize,>=Stealth,box/.style={draw,rounded corners=2pt,minimum height=0.7cm,inner sep=3pt,align=center}]
\node[box,fill=omDef!12] (py) at (0,0) {Appelant Python~:\\objets, tableaux numpy};
\node[box,fill=omMeth!15] (pb) at (4.6,0.9) {liaison pybind11~: vérifie\\le type et la disposition, relâche\\le verrou};
\node[box,fill=omThm!12] (ct) at (4.6,-0.9) {ctypes~: passe des pointeurs\\tels que déclarés, ne vérifie rien\\(l'enveloppe doit le faire)};
\node[box,fill=gray!10] (cpp) at (9.0,0.9) {Noyau C++20};
\node[box,fill=gray!10] (rs) at (9.0,-0.9) {Noyau Rust\\(ABI C, cdylib)};
\draw[->] (py) -- (pb); \draw[->] (py) -- (ct); \draw[->] (pb) -- (cpp); \draw[->] (ct) -- (rs);
\draw[dashed,omThm,rounded corners=3pt] (-1.6,-2.4) rectangle (10.4,-1.8);
\node[omThm] at (4.4,-2.1) {un seul tampon numpy~: lu et écrit sur place par les deux noyaux, jamais copié};
\end{tikzpicture}

\omcaption{Deux voies à travers la frontière. pybind11 génère une liaison qui vérifie le type et la disposition mémoire de chaque argument et peut
relâcher le verrou de l'interpréteur pendant que le C++ s'exécute~; ctypes charge n'importe quelle bibliothèque dotée d'une interface C et passe ce qu'on lui indique, si bien que l'enveloppe Python
doit faire la vérification. Les deux remettent au noyau un pointeur vers le même tampon numpy.}\label{fig:pl:native-extensions:boundary}
\end{omfigure}

La \Cref{fig:pl:native-extensions:calls} mesure la frontière. Une fonction à une étape appelée un million de fois depuis une boucle Python coûte 60~ns par appel
lorsqu'il s'agit d'une fonction Python, 118~ns avec pybind11 et 337~ns avec \texttt{ctypes}~: chaque appel natif convertit
trois flottants Python en doubles C et un en sens inverse, et \texttt{ctypes} le fait dynamiquement, d'après les types d'arguments déclarés à l'exécution.
L'arithmétique, une multiplication et une addition, prend moins d'une nanoseconde. Une fonction native appelée par élément est plus lente que le Python
qu'elle remplace, parce que la boucle qui l'entoure reste en Python.

\begin{omfigure}
\begin{tikzpicture}
\begin{axis}[width=11cm,height=4.4cm,xbar,bar width=9pt,xmin=0,xmax=480,xlabel={nanosecondes par appel, boucle comprise},y dir=reverse,
  ytick={0,1,2},yticklabels={fonction Python,C++ via pybind11,Rust via ctypes},
  tick label style={font=\small},label style={font=\small},grid=major,grid style={gray!20},point meta=explicit,
  nodes near coords,nodes near coords style={font=\footnotesize,anchor=west,/pgf/number format/fixed,/pgf/number format/precision=0},
  table/col sep=comma]
\addplot[fill=omThm!45,draw=omThm] table[x=ns_per_call,y expr=\coordindex,meta=ns_per_call]{figdata/platforms/09-native-extensions/measured_calls.csv};
\end{axis}
\end{tikzpicture}

\omcaption{Le coût d'un appel à une fonction à une étape, une fois par élément depuis une boucle Python, un million de fois~: comme fonction Python, comme C++
avec pybind11, et comme Rust avec \texttt{ctypes}. Mesuré sur un portable (Intel Core Ultra 7 155H) sous WSL2, un seul fil d'exécution, machine
par ailleurs inactive. Données~: \texttt{bench\_natext.py}.}\label{fig:pl:native-extensions:calls}
\end{omfigure}

\begin{proposition}[Déplacer la boucle, pas le corps]\label{prop:pl:native-extensions:loop}
Si un appel natif coûte $c_0$ par franchissement et $c_1$ par élément, et que le Python qu'il remplace coûte $p$ par élément, alors appeler la fonction
une fois par élément coûte $n(c_0 + c_1)$ et n'est plus rapide que Python que si $c_0 + c_1 < p$~; un appel unique sur l'ensemble du tableau coûte
$c_0 + n c_1$, ce qui est plus rapide que Python pour tout $n > c_0/(p - c_1)$ et tend vers $n c_1$.
\end{proposition}

\begin{proof}
Additionner les coûts des appels dans chaque arrangement et comparer avec $np$.
\end{proof}

Avec $c_0 \approx 118$~ns et $p \approx 60$~ns, le premier arrangement ne gagne jamais~; le second gagne au-delà de quelques éléments. Le
verrou de l'interpréteur est l'autre moitié de la frontière~: pendant que le noyau s'exécute, il ne touche à aucun objet Python, si bien que la liaison relâche le verrou
(voir le \cref{lst:pl:native-extensions:pybind}) et d'autres fils Python peuvent s'exécuter entre-temps (le tri du chapitre~8 faisait de même).

\section{Lier du C++ avec un lieur en-tête seul}

pybind11 est une bibliothèque C++ à en-tête seul~: la liaison est un fichier C++ qui l'inclut, décrit les arguments de chaque fonction, et est compilé
avec les en-têtes Python en une bibliothèque partagée que Python importe comme un module. \texttt{firm.natext} la compile au premier usage avec le
compilateur de la collection (\texttt{g++ -std=c++20 -O2 -shared -fPIC} avec les chemins d'inclusion donnés par \texttt{python3-config} et
pybind11).

\omcode{code/firm/natext/cpp/firm_natext_py.cpp}{11}{24}{La liaison pybind11 de la moyenne~: un tableau numpy pris comme tampon float64 contigu en ordre
C, un tableau de résultat alloué une seule fois, et le verrou de l'interpréteur relâché pendant que le C++ s'exécute.}\label{lst:pl:native-extensions:pybind}

\begin{definition}[Protocole de tampon]\label{def:pl:native-extensions:buffer}
Le \emph{protocole de tampon}\index{protocole de tampon} (buffer protocol) est l'interface de Python par laquelle un objet expose sa mémoire sous-jacente --- un pointeur,
le type des éléments, la forme et les pas --- à d'autres objets et au code compilé, afin qu'ils puissent la lire et l'écrire sans
la copier.
\end{definition}

Un tableau numpy expose sa mémoire de cette façon, et le \texttt{array\_t} de pybind11 la lit. La déclaration détermine ce qui se passe lorsque
le tableau n'est pas celui qu'attend le noyau. Avec \texttt{forcecast}, pybind11 convertit en silence~: un tableau float32 ou une vue à pas est copié
dans un nouveau tampon contigu float64, et une fonction censée écrire dans le tableau de l'appelant écrit à la place dans la copie. La
liaison de ce chapitre déclare \texttt{c\_style} sans \texttt{forcecast} et marque ses arguments \texttt{noconvert}, si bien que les deux sont refusés
par une erreur~: une copie doit être une décision, non une surprise.

\section{Lier Rust par l'ABI C}

Rust dispose aussi d'un générateur de liaisons pour Python~; cette série n'utilise aucune crate externe, si bien que le chapitre emprunte la voie qui n'en exige aucune~: une
bibliothèque de fonctions \texttt{extern \textquotedbl{}C\textquotedbl{}} compilées en bibliothèque dynamique compatible C, chargée avec le module \texttt{ctypes} de Python.

\omcode{code/firm/natext/rust/src/lib.rs}{29}{40}{Une fonction Rust à interface C~: pointeurs bruts et une longueur en entrée, une tranche construite à partir
d'eux sous des conditions énoncées, et le noyau sûr appelé sur les tranches.}\label{lst:pl:native-extensions:rust}

\begin{dated}{2026-09}{Liaisons Python de Rust}\label{dat:pl:native-extensions:pyo3}
PyO3, décrit dans son guide d'utilisation comme «~liaisons Rust pour Python, comprenant des outils pour créer des modules d'extension Python natifs~»
(consulté en septembre 2026), génère l'équivalent des vérifications et conversions d'arguments de pybind11 pour Rust. C'est une crate externe,
que la série ne construit pas~; l'interface C de ce chapitre est ce que produit une telle liaison par-dessous. Un \texttt{cdylib} est, dans
les termes de la Rust Reference, une bibliothèque système dynamique utilisée lorsqu'on compile une bibliothèque destinée à être chargée depuis un autre langage.
\end{dated}

\texttt{ctypes} est l'extrême opposé de pybind11~: il appelle toute fonction C avec les arguments qu'on lui indique et ne vérifie
rien de la mémoire située derrière un pointeur. L'enveloppe (\cref{lst:pl:native-extensions:ctypes}) déclare une fois pour toutes les types d'arguments de chaque fonction,
vérifie que chaque tableau est contigu en ordre C et du bon type, puis transmet les pointeurs~; sans ce contrôle, une vue à pas serait lue
comme si elle était contiguë, en silence.

\omcode{code/firm/natext/firm_natext.py}{105}{131}{Chargement de la bibliothèque Rust avec \texttt{ctypes}~: types d'arguments déclarés, et une
enveloppe qui vérifie ce que \texttt{ctypes} ne vérifie pas avant de transmettre des pointeurs vers les tampons numpy.}\label{lst:pl:native-extensions:ctypes}

\section{Sans copie à travers la frontière}

\begin{definition}[Transfert sans copie]\label{def:pl:native-extensions:zerocopy}
Un \emph{transfert sans copie} de données à travers une frontière\index{transfert sans copie} remet au destinataire une référence vers la mémoire de l'émetteur
(un pointeur et une disposition) plutôt qu'une copie, de sorte que son coût ne dépend pas de la taille des données~; il exige que les deux côtés s'accordent
sur la disposition, sur qui peut écrire et libérer la mémoire, et pendant combien de temps.
\end{definition}

Les deux liaisons sont sans copie pour leurs entrées~; \texttt{ewma\_into} l'est aussi pour sa sortie, en écrivant dans un tableau que
l'appelant a alloué (un test vérifie que l'adresse du tableau est inchangée). Ce sont les règles de propriété qui le rendent sûr~: le noyau lit
et n'écrit que pendant l'appel, ne conserve aucun pointeur après coup, et ne libère jamais ce que Python a alloué~; Python maintient les tableaux en vie pour
la durée de l'appel. Les \omterm{def:pl:a-tick-store-on-open-formats:mmap}{fichiers projetés en mémoire} du chapitre~4 étendent la même idée au disque et à d'autres processus, et la
disposition d'Arrow du chapitre~3 à d'autres bibliothèques et langages.

\begin{omfigure}
\begin{tikzpicture}
\begin{axis}[width=12cm,height=6cm,xmode=log,ymode=log,xlabel={longueur du tableau (échelle log)},ylabel={ns par élément (échelle log)},
  tick label style={font=\small},label style={font=\small},grid=major,grid style={gray!20},
  legend columns=4,legend style={font=\footnotesize,at={(0.5,-0.25)},anchor=north,draw=none,/tikz/every even column/.append style={column sep=6pt}},
  table/col sep=comma]
\addplot[thick,mark=*,omDef] table[x=n,y=ewma_numpy]{figdata/platforms/09-native-extensions/measured_sweep.csv};
\addplot[thick,mark=square*,omMeth] table[x=n,y=ewma_cpp]{figdata/platforms/09-native-extensions/measured_sweep.csv};
\addplot[thick,mark=triangle*,omThm] table[x=n,y=ewma_rust]{figdata/platforms/09-native-extensions/measured_sweep.csv};
\addplot[thick,mark=diamond*,gray] table[x=n,y=ewma_python]{figdata/platforms/09-native-extensions/measured_sweep.csv};
\legend{numpy (filtre),C++ (pybind11),Rust (ctypes),Python pur}
\end{axis}
\end{tikzpicture}

\omcaption{Temps par élément de la moyenne sur tout le tableau, en fonction de la longueur de celui-ci. Pour les tableaux courts, chaque voie est dominée par son
coût fixe par appel ($c_0/n$)~; pour les tableaux longs, les boucles natives se stabilisent vers 1.3 à 1.9~ns par élément et le filtre de numpy vers 3 à 3.6, tandis que
Python pur reste à des centaines. Mesures faites sur un ordinateur portable (Intel Core Ultra 7 155H) sous WSL2, un seul fil d'exécution, machine par ailleurs au repos. Données~:
\texttt{bench\_natext.py}.}\label{fig:pl:native-extensions:sweep}
\end{omfigure}

Les \Cref{fig:pl:native-extensions:sweep,fig:pl:native-extensions:asof} sont les variantes sur tout le tableau. Pour la moyenne, les deux voies natives
battent le filtre de numpy à toutes les longueurs mesurées~: la routine de filtrage a un coût fixe de plusieurs microsecondes, supérieur à celui de pybind11~; à
un million d'éléments, C++ prend 1.4~ns par élément contre 3.3 pour numpy. Pour la recherche as-of, c'est l'algorithme qui compte davantage que le
langage~: la recherche dichotomique de numpy, élément par élément, coûte davantage à mesure que les tableaux grossissent, 21~ns à un million, tandis que le parcours de fusion reste à
environ 6, soit trois fois moins. La voie \texttt{ctypes} paie un coût fixe plus élevé (les conversions de pointeurs de l'enveloppe) et
ne dépasse numpy qu'à partir d'environ mille éléments pour la recherche.

\begin{omfigure}
\begin{tikzpicture}
\begin{axis}[width=12cm,height=5.6cm,xmode=log,ymode=log,xlabel={longueur du tableau (échelle log)},ylabel={ns par élément (échelle log)},
  tick label style={font=\small},label style={font=\small},grid=major,grid style={gray!20},
  legend columns=3,legend style={font=\footnotesize,at={(0.5,-0.25)},anchor=north,draw=none,/tikz/every even column/.append style={column sep=6pt}},
  table/col sep=comma]
\addplot[thick,mark=*,omDef] table[x=n,y=asof_numpy]{figdata/platforms/09-native-extensions/measured_sweep.csv};
\addplot[thick,mark=square*,omMeth] table[x=n,y=asof_cpp]{figdata/platforms/09-native-extensions/measured_sweep.csv};
\addplot[thick,mark=triangle*,omThm] table[x=n,y=asof_rust]{figdata/platforms/09-native-extensions/measured_sweep.csv};
\legend{numpy (recherche dichotomique par élément),C++ (fusion),Rust (fusion)}
\end{axis}
\end{tikzpicture}

\omcaption{Temps par élément de la recherche as-of de $n$ instants triés dans $n$ autres. La recherche dichotomique de numpy coûte davantage par élément à mesure que les
tableaux grossissent~; le parcours de fusion natif coûte un petit nombre constant de nanosecondes une fois le coût fixe amorti. Mesures faites sur un ordinateur portable (Intel Core
Ultra 7 155H) sous WSL2, un seul fil d'exécution, machine par ailleurs au repos. Données~: \texttt{bench\_natext.py}.}\label{fig:pl:native-extensions:asof}
\end{omfigure}

\begin{method}[Écrire un noyau natif pour la plateforme de recherche]\label{met:pl:native-extensions:method}
(1) Écrire le noyau en C++ ou en Rust sans rien connaître de Python~: pointeurs, longueurs, types simples. (2) Le tester isolément, dans son propre langage.
(3) Lier le noyau une fois par tableau, jamais une fois par élément~: des tableaux en entrée, des tableaux en sortie, la boucle à l'intérieur. (4) Refuser les dispositions et les types que le noyau
ne lit pas sur place, plutôt que de les copier en silence. (5) Relâcher le verrou de l'interpréteur autour du calcul pur. (6) Tester la liaison
contre la référence Python sur les mêmes données, y compris les entrées vides et les cas limites.
\end{method}

\section{Tutoriel~: un noyau, quatre implémentations}\label{tut:pl:native-extensions:tut}

\begin{tutorial}
\textbf{Objectif.} Construire les noyaux C++ et Rust, les lier, montrer que quatre implémentations concordent, et mesurer la frontière.
\textbf{État final~:} les \cref{fig:pl:native-extensions:calls,fig:pl:native-extensions:sweep,fig:pl:native-extensions:asof}.
\begin{tutsteps}
\item \textbf{Noyaux}~: lire le \cref{lst:pl:native-extensions:kernels}~; compiler et exécuter \texttt{firm\_natext\_test.cpp}~; exécuter
\texttt{cargo test} dans \texttt{code/firm/natext/rust}.
\item \textbf{Liaison}~: \texttt{firm\_natext.cpp\_module()} compile et importe le module pybind11~;
\texttt{rust\_library()} construit le \texttt{cdylib} et déclare ses fonctions à \texttt{ctypes}.
\item \textbf{Concordance}~: les tests d'acceptation comparent les quatre moyennes et les quatre recherches, sur des entrées aléatoires et limites.
\item \textbf{Refus}~: passer une vue à pas et un tableau float32 à chaque liaison, et lire les erreurs.
\item \textbf{Mesure}~: \texttt{bench\_natext.py} (coût par appel et balayage des tailles).
\end{tutsteps}
\textbf{Pour aller plus loin.} Ajouter \texttt{forcecast} à la liaison et observer \texttt{ewma\_into} écrire dans une copie~; appeler le noyau Rust
depuis deux fils d'exécution Python, via \texttt{ctypes} (qui relâche le verrou autour des appels étrangers), et mesurer le gain de vitesse.
\end{tutorial}

\section{Construction~: noyaux natifs pour la plateforme de recherche}\label{bld:pl:native-extensions:natext}

\begin{build}
\bfield{Rôle} Le motif et la mécanique de compilation que suit chaque noyau natif de la plateforme~: les lecteurs de la \omterm{def:pl:a-tick-store-on-open-formats:store}{base de ticks} du
chapitre~4, les boucles internes du moteur de backtest (chapitre~11) et l'évaluation du compilateur de paiements (chapitre~19) peuvent
tous franchir la frontière de cette façon.
\bfield{Interface} \texttt{ewma\_python}, \texttt{ewma\_numpy}, \texttt{ewma\_cpp}, \texttt{ewma\_cpp\_into}, \texttt{ewma\_rust}~;
\texttt{asof\_python}, \texttt{asof\_numpy}, \texttt{asof\_cpp}, \texttt{asof\_rust}~; \texttt{step\_python}, \texttt{step\_cpp},
\texttt{step\_rust}~; \texttt{cpp\_module()}, \texttt{rust\_library()}. En-tête C++20 \texttt{firm::natext} (\texttt{ewma},
\texttt{ewma\_step}, \texttt{asof\_index})~; crate Rust \texttt{natext\_rs} avec \texttt{natext\_ewma}, \texttt{natext\_asof},
\texttt{natext\_ewma\_step}.
\bfield{Règles} Les noyaux ne connaissent rien de Python~; un appel par tableau~; les entrées sont lues sur place ou refusées~; le verrou
est relâché autour du calcul~; aucun pointeur n'est conservé après un appel~; les compilations font partie des tests, et une chaîne d'outils absente les fait échouer.
\bfield{Tests d'acceptation} \texttt{code/firm/natext/tests/}, \texttt{cpp/firm\_natext\_test.cpp}, \texttt{rust/}~: quatre moyennes et quatre
recherches égales, bornes comprises~; trois pas égaux~; sortie écrite sur place~; tableaux à pas et mal typés refusés par les deux liaisons.
\bfield{Pour aller plus loin} Un noyau sur les enregistrements plats de la \omterm{def:pl:a-tick-store-on-open-formats:store}{base de ticks} (un tampon structuré, chapitre~4)~; un filtre du
second ordre dans les quatre langages~; les liaisons comparées sous deux fils d'exécution.
\end{build}

\begin{omsources}
\item Documentation de pybind11, \emph{NumPy} (protocole de tampon, \texttt{array\_t}, \texttt{c\_style}, \texttt{forcecast},
\texttt{noconvert}).
\item T.~Oliphant et C.~Banks, \emph{PEP 3118 -- Revising the buffer protocol} (Final).
\item The Rust Reference, \emph{Linkage} (\texttt{cdylib})~; PyO3 user guide.
\end{omsources}

\section{Exercices}

\begin{exercise}[$\star$]\label{exo:pl:native-extensions:1}
À 118~ns par appel via pybind11 et 60~ns par pas en Python, combien de temps prend une journée de deux cents millions de pas dans chaque cas,
appelée pour chaque élément~? Et calculée par un seul appel à 1.4~ns par élément~?
\end{exercise}

\begin{exercise}[$\star$]\label{exo:pl:native-extensions:2}
Pourquoi l'appel \texttt{ctypes} est-il plus lent que l'appel pybind11, alors que les deux se terminent par les mêmes instructions machine~?
\end{exercise}

\begin{exercise}[$\star$]\label{exo:pl:native-extensions:3}
Que fait \texttt{forcecast} d'un tableau float32 passé à une fonction déclarée pour float64, et pourquoi est-ce dangereux pour
\texttt{ewma\_into}~?
\end{exercise}

\begin{exercise}[$\star\star$]\label{exo:pl:native-extensions:4}
Avec la \cref{prop:pl:native-extensions:loop}, $c_0 = 118$~ns, $c_1 = 1.4$~ns et $p = 60$~ns, à partir de combien d'éléments un appel sur le
tableau entier est-il plus rapide que la boucle Python~?
\end{exercise}

\begin{exercise}[$\star\star$]\label{exo:pl:native-extensions:5}
Pourquoi la recherche as-of de numpy coûte-t-elle davantage par élément à mesure que les tableaux grandissent, alors que la passe de fusion ne le fait pas~?
\end{exercise}

\begin{exercise}[$\star\star$]\label{exo:pl:native-extensions:6}
Un noyau conserve un pointeur vers un tableau numpy entre deux appels pour le réutiliser. Décrire comment cela tourne mal, et ce que la règle de propriété interdit.
\end{exercise}

\begin{exercise}[$\star\star\star$]\label{exo:pl:native-extensions:7}
\emph{Programmation.} Avec les fonctions du chapitre, calculer la recherche as-of de 100\,000 temps triés aléatoires dans 50\,000 autres par les quatre
implémentations, et vérifier qu'elles concordent, y compris pour un temps de gauche avant le premier temps de droite et un temps après le dernier.
\end{exercise}

\begin{exercise}[$\star\star\star$]\label{exo:pl:native-extensions:8}
\emph{Trouver la faille.} «~Notre extension C++ prend une liste de flottants Python, la convertit en \texttt{std::vector<double>}, calcule, et
renvoie une nouvelle liste. Elle est plus rapide que numpy sur notre banc d'essai de dix nombres, donc nous l'utilisons partout.~»
\end{exercise}

\section{Problème~: le million d'appels}

\begin{problem}[{Problème du week-end --- où franchir la frontière}]\label{pb:pl:native-extensions:1}
Les noyaux du chapitre, leurs quatre implémentations, les deux liaisons et les mesures de \texttt{bench\_natext.py}.

\medskip\noindent\textbf{Partie I --- La traversée.}
\begin{enumerate}
\item Que fait un appel natif en plus du calcul arithmétique~?
\item Quels sont les coûts mesurés par appel pour les trois voies~?
\item Pourquoi un appel natif par élément est-il plus lent que la fonction Python qu'il remplace~?
\item Que permet le relâchement du verrou de l'interpréteur, et quand est-il sûr~?
\item Qu'est-ce qu'une \omterm{def:pl:native-extensions:ffi}{interface binaire d'application}, et pourquoi l'interface C est-elle le terrain commun~?
\end{enumerate}

\medskip\noindent\textbf{Partie II --- Les deux liaisons.}
\begin{enumerate}[resume]
\item Que vérifie pybind11 que \texttt{ctypes} ne vérifie pas~?
\item Que doit vérifier le wrapper \texttt{ctypes} lui-même~?
\item Que se passe-t-il pour une vue à pas passée à chaque liaison~?
\item Pourquoi \texttt{forcecast} n'est-il pas utilisé~?
\item Qu'ajouterait PyO3 à la voie Rust~?
\end{enumerate}

\medskip\noindent\textbf{Partie III --- Tableaux.}
\begin{enumerate}[resume]
\item Quels sont les coûts par élément de la moyenne sur tout le tableau à un million d'éléments, par voie~?
\item Pourquoi les voies natives battent-elles le filtre de numpy à chaque longueur mesurée~?
\item Quels sont les coûts par élément de la recherche as-of à un million, par voie~?
\item À partir de quelle longueur la recherche Rust bat-elle celle de numpy~?
\item Lequel des deux noyaux gagne le plus à passer en natif, et pourquoi~?
\end{enumerate}

\medskip\noindent\textbf{Partie IV --- Le verdict.}
\begin{enumerate}[resume]
\item Énoncer le \emph{résultat nommé}~: le coût par appel de chaque liaison, les longueurs au-delà desquelles chaque version native bat numpy, et le
gain de déplacer la boucle, et non le corps, au franchissement de la frontière.
\item Combien de temps une journée de deux cents millions de pas économise-t-elle en déplaçant la boucle en C++~?
\item Quand une \omterm{def:pl:native-extensions:extension}{extension native} ne vaut-elle pas la peine d'être écrite~?
\item Qu'exige la règle de propriété du noyau et de l'appelant~?
\item En une phrase~: où doit se trouver la frontière entre Python et le code natif~?
\end{enumerate}
\end{problem}

\section{Questions d'entretien}

\begin{interviewq}[$\star$ \iqroles{développeur, chercheur}]\label{iq:pl:native-extensions:1}
Vous avez écrit une fonction en C++ pour accélérer une boucle Python, et elle est devenue plus lente~? Que s'est-il passé~?
\end{interviewq}

\begin{interviewq}[$\star\star$ \iqroles{développeur}]\label{iq:pl:native-extensions:2}
Comment passez-vous un tableau numpy à C++ ou à Rust sans le copier, et que peut-il mal se passer~?
\end{interviewq}

\begin{interviewq}[$\star\star$ \iqroles{développeur}]\label{iq:pl:native-extensions:3}
Comparez pybind11, ctypes et un générateur de liaisons Rust pour exposer un noyau numérique à Python.
\end{interviewq}

\begin{interviewq}[$\star\star$ \iqroles{développeur}]\label{iq:pl:native-extensions:4}
Pourquoi une \omterm{def:pl:native-extensions:extension}{extension native} doit-elle relâcher le verrou de l'interpréteur, et que ne doit-elle pas faire pendant qu'il est relâché~?
\end{interviewq}

\begin{interviewq}[$\star\star\star$ \iqroles{développeur, chercheur}]\label{iq:pl:native-extensions:5}
Concevez la couche native d'une plateforme de recherche~: quels noyaux passent en natif, comment ils sont liés, construits, testés et distribués.
\end{interviewq}
